\documentclass[12pt,a4paper]{article}

% ------------------------------------------------
% PAQUETES
% ------------------------------------------------

\usepackage[utf8]{inputenc}
\usepackage[T1]{fontenc}
\usepackage[spanish,es-nodecimaldot]{babel}

\usepackage{amsmath}
\usepackage{amssymb}
\usepackage{amsthm}
\usepackage{mathtools}

\usepackage{graphicx}
\usepackage{float}
\usepackage{booktabs}
\usepackage{array}
\usepackage{enumitem}

\usepackage{algorithm}
\usepackage{algpseudocode}

\usepackage[round,authoryear]{natbib}
\usepackage[hidelinks]{hyperref}

\usepackage{geometry}
\geometry{
    top=2.5cm,
    bottom=2.5cm,
    left=2.8cm,
    right=2.8cm
}

\usepackage{setspace}
\onehalfspacing

\usepackage{parskip}

\usepackage{fancyhdr}
\pagestyle{fancy}
\fancyhf{}
\rhead{Recocido simulado y el TSP}
\lhead{Cadenas de Markov}
\cfoot{\thepage}

\graphicspath{{../results/}}

% ------------------------------------------------
% ENTORNOS
% ------------------------------------------------

\theoremstyle{definition}
\newtheorem{definicion}{Definición}[section]
\newtheorem{ejemplo}[definicion]{Ejemplo}
\newtheorem{observacion}[definicion]{Observación}

\theoremstyle{plain}
\newtheorem{teorema}[definicion]{Teorema}
\newtheorem{proposicion}[definicion]{Proposición}
\newtheorem{lema}[definicion]{Lema}
\newtheorem{corolario}[definicion]{Corolario}

\renewcommand{\proofname}{Demostración}

% ------------------------------------------------
% NOTACIÓN
% ------------------------------------------------

\newcommand{\R}{\mathbb{R}}
\newcommand{\Prob}{\mathbb{P}}
\newcommand{\E}{\mathbb{E}}
\newcommand{\Var}{\operatorname{Var}}
\newcommand{\TV}{\mathrm{TV}}
\newcommand{\Sopt}{S^{\ast}}
\newcommand{\cmax}{c_{\max}}

% ------------------------------------------------
% INFORMACIÓN DEL DOCUMENTO
% ------------------------------------------------

\title{
    \textbf{El recocido simulado como cadena de Markov: una aplicación al problema del agente viajero}
}

\author{
    Henry Giovani Barrera Martínez
    Laura Victoria Torres Pinilla
}

\date{\today}

% ------------------------------------------------
% DOCUMENTO
% ------------------------------------------------

\begin{document}

\begin{titlepage}

    \centering

    \vspace*{2cm}

    {\Large \textbf{Universidad Nacional de Colombia}\par}

    \vspace{0.5cm}

    {\large Departamento de Matemáticas\par}

    \vspace{2cm}

    {\LARGE
    \textbf{El recocido simulado como cadena de Markov: una aplicación al problema del agente viajero}
    \par}

    \vspace{1.5cm}

    {\large Asignatura: Cadenas de Markov\par}

    {\large Proyecto final\par}

    \vspace{2cm}

    {\large
    \textbf{Estudiantes:}\\
        Henry Giovani Barrera Martínez\\
        Laura Victoria Torres Pinilla
    \par}

    \vfill

    {\large Bogotá D.C., Colombia\par}

    {\large \today\par}

\end{titlepage}

% ------------------------------------------------
% RESUMEN
% ------------------------------------------------

\begin{abstract}
El recocido simulado es una de las metaheurísticas más usadas para el
problema del agente viajero (TSP) y, al mismo tiempo, una cadena de
Markov cuyo comportamiento puede analizarse con rigor. En este trabajo
se construye la cadena de Metropolis sobre el espacio de rutas con la
vecindad 2-opt y se demuestra que es irreducible, aperiódica y
reversible respecto de la distribución de Boltzmann. Se estudia además
cómo la temperatura controla a la vez la concentración de esa
distribución en las rutas óptimas y la velocidad de convergencia, medida
por la brecha espectral. Con ello, el recocido aparece como una cadena
no homogénea cuya convergencia caracteriza el teorema de Hajek. También
se analiza una variante sobre la formulación QUBO del TSP y se demuestra
que la penalización de las restricciones crea barreras de energía que
crecen linealmente con el parámetro \(M\).

Los experimentos, reproducibles en un notebook de Google Colab,
confirman la teoría. Con cinco ciudades, las frecuencias empíricas
convergen a la distribución exacta y la pendiente de \(-\ln\gamma(T)\)
frente a \(1/T\) es \(68.2\), a un \(3\,\%\) de la profundidad crítica
exacta \(h^\ast=66\). El recocido sobre rutas encuentra el óptimo en
las \(40\) corridas sobre cuatro instancias de TSPLIB, hasta
\texttt{berlin52}. En cambio, sobre el QUBO el desempeño depende de
\(M\) en forma de U y, en \texttt{burma14}, no baja del \(29\,\%\) por
encima del óptimo. La conclusión es que, para una misma cadena de
Metropolis, la elección del espacio de estados determina la geometría
del paisaje y, con ella, la eficacia del algoritmo.

\medskip
\noindent\textbf{Palabras clave:} cadenas de Markov, algoritmo de
Metropolis, recocido simulado, brecha espectral, problema del agente
viajero, QUBO.
\end{abstract}


\tableofcontents

\newpage

% ------------------------------------------------
% 1. INTRODUCCIÓN
% ------------------------------------------------

\section{Introducción}
\label{sec:introduccion}

TODO

% ------------------------------------------------
% 2. DESCRIPCIÓN DEL PROBLEMA Y OBJETIVOS
% ------------------------------------------------

\section{Descripción del problema y objetivos}
\label{sec:problema}

\subsection{El problema del agente viajero}
\label{subsec:tsp}

Se tienen \(n\) ciudades \(V=\{1,\ldots,n\}\) y un costo
\(c_{ij}=c_{ji}\ge0\) de viajar entre cada par de ciudades distintas;
en las instancias de este trabajo, \(c_{ij}\) es una distancia. El TSP
simétrico consiste en encontrar una permutación cíclica de las ciudades,
es decir, una \emph{ruta}, cuyo costo total
\[
f(\sigma)=\sum_{k=1}^{n}c_{\sigma_k\sigma_{k+1}},\qquad
\sigma_{n+1}=\sigma_1,
\]
sea mínimo. Como el costo no depende de la ciudad de partida ni del
sentido del recorrido, el número de rutas distintas es
\((n-1)!/2\). La tabla~\ref{tab:explosion} muestra este número para los
tamaños usados en este trabajo.

\begin{table}[H]
\centering
\small
\begin{tabular}{@{}lrr@{}}
\toprule
Instancia & \(n\) & rutas \((n-1)!/2\) \\
\midrule
aleatoria          &  5 & \(12\) \\
aleatoria          &  8 & \(2\,520\) \\
\texttt{burma14}   & 14 & \(3.1\times10^{9}\) \\
\texttt{ulysses16} & 16 & \(6.5\times10^{11}\) \\
\texttt{ulysses22} & 22 & \(2.6\times10^{19}\) \\
\texttt{berlin52}  & 52 & \(7.8\times10^{65}\) \\
\bottomrule
\end{tabular}
\caption{Número de rutas distintas según el número de ciudades.}
\label{tab:explosion}
\end{table}

Con \(n=5\) es posible enumerar todas las rutas, y por eso esa instancia
se usa para calcular la cadena de Markov de forma exacta. Ya con
\(n=22\), en cambio, evaluar una ruta por nanosegundo tomaría más de
ochocientos años, y con \(n=52\) el número de rutas supera con creces
el número estimado de átomos en la Tierra. La exploración exhaustiva es
inviable, y aquí entran los métodos aleatorizados: el recocido simulado
no recorre el espacio de rutas, sino que lo muestrea mediante una
cadena de Markov diseñada para concentrarse en las rutas de bajo costo.

\subsection{Objetivos}
\label{subsec:objetivos}

\subsubsection*{Objetivo general}

Estudiar el recocido simulado aplicado al problema del agente viajero
como una cadena de Markov, demostrando sus propiedades fundamentales y
contrastándolas con simulaciones reproducibles.

\subsubsection*{Objetivos específicos}

\begin{itemize}
    \item Construir la cadena de Metropolis sobre el espacio de rutas con
    la vecindad 2-opt y demostrar que es irreducible, aperiódica y
    reversible respecto de la distribución de Boltzmann.
    \item Analizar el efecto de la temperatura sobre la distribución
    estacionaria y sobre la velocidad de convergencia, medida por la
    brecha espectral, y relacionar esta última con la profundidad
    crítica del paisaje de costos.
    \item Presentar el recocido simulado como una cadena de Markov no
    homogénea y discutir, con el teorema de Hajek, las diferencias entre
    el enfriamiento logarítmico y el geométrico.
    \item Formular la variante del recocido sobre el modelo QUBO del TSP
    y explicar, a partir de la geometría del paisaje de energía, cómo
    influye el parámetro de penalización \(M\).
    \item Implementar la cadena exacta para \(n=5\) y el recocido
    simulado en un notebook reproducible, y verificar con él los
    resultados teóricos en instancias aleatorias y de TSPLIB.
\end{itemize}

% ------------------------------------------------
% 3. MARCO TEÓRICO Y FORMULACIÓN MATEMÁTICA DEL MODELO
% ------------------------------------------------

\section{Marco teórico y formulación matemática del modelo}
\label{sec:marco}

Esta sección construye la cadena de Markov que subyace al recocido
simulado aplicado al problema del agente viajero (TSP). Primero se
recuerdan los resultados de cadenas de Markov finitas que se usarán
(\S\ref{subsec:preliminares}). Luego se definen el espacio de estados y
la vecindad 2-opt (\S\ref{subsec:estados}), se demuestra la
irreducibilidad de la vecindad (\S\ref{subsec:irreducibilidad}) y se
construye la cadena de Metropolis a temperatura fija, cuya distribución
estacionaria es la de Boltzmann (\S\ref{subsec:metropolis}). Después se
estudian el papel de la temperatura (\S\ref{subsec:temperatura}) y la
velocidad de convergencia a través de la brecha espectral
(\S\ref{subsec:brecha}). Con esto, el recocido simulado aparece como
una cadena no homogénea cuya convergencia caracteriza el teorema de
\citet{hajek1988} (\S\ref{subsec:recocido}). Finalmente se formula la
variante sobre el modelo QUBO y se muestra por qué su paisaje de
energía es más difícil de recorrer (\S\ref{subsec:qubo}).

% ------------------------------------------------
\subsection{Preliminares: cadenas de Markov finitas}
\label{subsec:preliminares}

Sea \(S\) un conjunto finito y \(P=(P(x,y))_{x,y\in S}\) una matriz
estocástica, es decir, \(P(x,y)\ge 0\) y \(\sum_{y\in S}P(x,y)=1\) para
todo \(x\in S\). Una sucesión de variables aleatorias
\((X_t)_{t\ge 0}\) con valores en \(S\) es una \emph{cadena de Markov
homogénea} con matriz de transición \(P\) si
\[
\Prob(X_{t+1}=y \mid X_t=x,\, X_{t-1}=x_{t-1},\ldots,X_0=x_0)
=P(x,y)
\]
para todo \(t\ge 0\) y toda historia con probabilidad positiva. Se
escribe \(P^t(x,y)=\Prob(X_t=y\mid X_0=x)\).

\begin{definicion}[Irreducibilidad, período y reversibilidad]
\label{def:propiedades}
Sea \(P\) una matriz estocástica sobre \(S\) finito.
\begin{enumerate}[label=(\roman*)]
    \item \(P\) es \emph{irreducible} si para todo \(x,y\in S\) existe
    \(t\ge 1\) con \(P^t(x,y)>0\).
    \item El \emph{período} de \(x\) es
    \(d(x)=\gcd\{t\ge 1: P^t(x,x)>0\}\). Si \(P\) es irreducible,
    todos los estados tienen el mismo período, y \(P\) es
    \emph{aperiódica} cuando ese período es \(1\).
    \item Una distribución \(\pi\) sobre \(S\) es \emph{estacionaria}
    si \(\pi P=\pi\).
    \item \(P\) es \emph{reversible} respecto de \(\pi\) si satisface
    las \emph{ecuaciones de balance detallado}
    \begin{equation}
    \label{eq:balance}
    \pi(x)P(x,y)=\pi(y)P(y,x)\qquad\text{para todo } x,y\in S.
    \end{equation}
\end{enumerate}
\end{definicion}

Si \(P\) es irreducible, basta que un estado \(x\) tenga
\(P(x,x)>0\) para que la cadena sea aperiódica, pues entonces
\(1\in\{t: P^t(x,x)>0\}\). Además, sumando \eqref{eq:balance} sobre
\(x\) se obtiene \(\sum_x\pi(x)P(x,y)=\pi(y)\), es decir, \emph{toda
distribución que satisface el balance detallado es estacionaria}.

Los tres resultados siguientes son clásicos y se usan sin
demostración; pueden consultarse en
\citet[caps.~1 y~4]{levin2017} y en \citet[caps.~3 y~4]{bremaud2020}.

\begin{teorema}[Existencia y unicidad]
\label{teo:unicidad}
Si \(P\) es irreducible sobre \(S\) finito, existe una única
distribución estacionaria \(\pi\), y \(\pi(x)>0\) para todo
\(x\in S\).
\end{teorema}

\begin{teorema}[Convergencia]
\label{teo:convergencia}
Si \(P\) es irreducible y aperiódica sobre \(S\) finito, con
distribución estacionaria \(\pi\), entonces existen constantes
\(C>0\) y \(\alpha\in(0,1)\) tales que
\[
\max_{x\in S}\,\bigl\|P^t(x,\cdot)-\pi\bigr\|_{\TV}\le C\alpha^t
\qquad\text{para todo } t\ge 0,
\]
donde
\(\|\mu-\nu\|_{\TV}=\tfrac12\sum_{y\in S}|\mu(y)-\nu(y)|\).
\end{teorema}

\begin{teorema}[Teorema ergódico]
\label{teo:ergodico}
Si \(P\) es irreducible sobre \(S\) finito, con distribución
estacionaria \(\pi\), entonces para toda función \(g:S\to\R\) y toda
distribución inicial,
\[
\frac1t\sum_{s=0}^{t-1}g(X_s)\xrightarrow[t\to\infty]{\text{c.s.}}
\sum_{x\in S}g(x)\pi(x).
\]
En particular, tomando \(g=\mathbf 1_{\{y\}}\), la fracción de tiempo
que la cadena pasa en \(y\) converge casi seguramente a \(\pi(y)\).
\end{teorema}

El teorema~\ref{teo:ergodico} es la base del primer experimento de la
sección~\ref{sec:resultados}: las frecuencias empíricas de visita de
una trayectoria larga se comparan con la distribución estacionaria
calculada de forma exacta.

% ------------------------------------------------
\subsection{Espacio de estados y vecindad 2-opt}
\label{subsec:estados}

Sea \(V=\{1,\ldots,n\}\) un conjunto de \(n\ge 4\) ciudades con costos
simétricos \(c_{ij}=c_{ji}\ge 0\) para \(i\ne j\), y sea
\(\cmax=\max_{i\ne j}c_{ij}\). Una \emph{ruta} es un ciclo
hamiltoniano del grafo completo \(K_n\), es decir, un conjunto
\(x\) de \(n\) aristas que visita cada ciudad exactamente una vez. Su
costo es
\[
f(x)=\sum_{\{i,j\}\in x}c_{ij}.
\]
El espacio de estados de la cadena es el conjunto \(S_n\) de todas las
rutas. Como los costos son simétricos, una ruta no depende del punto de
partida ni del sentido de recorrido, de modo que
\begin{equation}
\label{eq:num_rutas}
|S_n|=\frac{(n-1)!}{2}.
\end{equation}
Para \(n=5\) se tiene \(|S_5|=12\), lo que permite calcular la cadena de
forma exacta. Para \(n=52\) (la instancia \texttt{berlin52} de TSPLIB),
\(|S_{52}|\approx 7.8\times10^{65}\). Denotamos por
\(f^\ast=\min_{x\in S_n}f(x)\) el costo óptimo y por
\(\Sopt=\{x\in S_n:f(x)=f^\ast\}\) el conjunto de rutas óptimas.

Es útil representar una ruta como una sucesión
\((\sigma_1,\sigma_2,\ldots,\sigma_n)\) con \(\sigma_1=1\), entendida de
forma cíclica (\(\sigma_{n+1}=\sigma_1\)). Cada ruta admite exactamente
dos representaciones de este tipo, una por cada sentido de recorrido.

\begin{definicion}[Movimiento 2-opt]
\label{def:2opt}
Sea \(x\in S_n\). Un \emph{movimiento 2-opt} retira de \(x\) dos aristas
no adyacentes \(\{a,b\}\) y \(\{c,d\}\), donde \(a,b,c,d\) aparecen en
ese orden a lo largo de la ruta, y las reemplaza por \(\{a,c\}\) y
\(\{b,d\}\) \citep{croes1958}. En la representación como sucesión, esto
equivale a invertir el segmento comprendido entre \(b\) y \(c\):
\[
(\ldots,a,\underbrace{b,\ldots,c}_{\text{se invierte}},d,\ldots)
\;\longmapsto\;
(\ldots,a,c,\ldots,b,d,\ldots).
\]
La \emph{vecindad} \(N(x)\) es el conjunto de rutas que se obtienen de
\(x\) mediante un movimiento 2-opt.
\end{definicion}

La otra forma de reconectar, \(\{a,d\}\) y \(\{b,c\}\), divide la ruta en
dos subciclos, de modo que la reconexión del movimiento 2-opt es la
única que produce otra ruta.

\begin{proposicion}[Estructura de la vecindad]
\label{prop:vecindad}
Para \(n\ge 4\) y toda \(x\in S_n\):
\begin{enumerate}[label=(\alph*)]
    \item \(|N(x)|=\dfrac{n(n-3)}{2}\);
    \item \(y\in N(x)\) si y solo si \(x\in N(y)\).
\end{enumerate}
\end{proposicion}

\begin{proof}
(a) Una arista de un ciclo de \(n\) aristas es adyacente a exactamente
dos aristas. Por lo tanto, cada arista forma par no adyacente con
\(n-3\) aristas, y el número de pares no adyacentes es \(n(n-3)/2\).
Cada par da una ruta \(y\ne x\), porque \(y\) contiene las aristas
\(\{a,c\}\) y \(\{b,d\}\), que no están en \(x\). Además, pares distintos
dan rutas distintas, porque el par retirado se recupera como
\(x\setminus y\).

(b) Si \(y\) se obtiene de \(x\) retirando \(\{a,b\},\{c,d\}\) y
agregando \(\{a,c\},\{b,d\}\), entonces estas dos últimas aristas no
comparten vértices, luego son no adyacentes en \(y\), y retirarlas y
reconectar devuelve \(x\).
\end{proof}

La propuesta que se implementa elige un par de posiciones \(\{i,j\}\)
uniformemente entre los \(\binom n2\) posibles y aplica el movimiento
2-opt que retira las aristas que salen de las posiciones \(i\) y \(j\).
Los \(n\) pares que corresponden a aristas adyacentes no cambian la ruta
(\emph{movimientos nulos}), y por la proposición~\ref{prop:vecindad}
cada uno de los \(n(n-3)/2\) pares restantes produce un vecino distinto.
La \emph{matriz de propuesta} es entonces
\begin{equation}
\label{eq:propuesta}
R(x,y)=
\begin{cases}
\dfrac{2}{n(n-1)}, & y\in N(x),\\[1.5ex]
\dfrac{2}{n-1}, & y=x,\\[1ex]
0, & \text{en otro caso},
\end{cases}
\qquad R(x,y)=R(y,x).
\end{equation}
Es simétrica por la proposición~\ref{prop:vecindad}(b), y la
probabilidad positiva de quedarse quieto tendrá un papel en la
aperiodicidad y en el espectro.

% ------------------------------------------------
\subsection{Irreducibilidad de la vecindad 2-opt}
\label{subsec:irreducibilidad}

\begin{proposicion}
\label{prop:irreducible}
Para \(n\ge 4\), la matriz de propuesta \(R\) es irreducible. Más aún,
desde cualquier ruta se llega a cualquier otra con a lo sumo \(n-2\)
movimientos 2-opt.
\end{proposicion}

\begin{proof}
Representemos las rutas como sucesiones con \(\sigma_1=1\). Para
\(2\le i<j\le n\), sea \(\rho_{ij}\) la inversión del segmento
\((\sigma_i,\ldots,\sigma_j)\). Esta inversión retira las aristas
\(\{\sigma_{i-1},\sigma_i\}\) y \(\{\sigma_j,\sigma_{j+1}\}\) y agrega
\(\{\sigma_{i-1},\sigma_j\}\) y \(\{\sigma_i,\sigma_{j+1}\}\). Las aristas
retiradas son no adyacentes salvo cuando \(\sigma_{i-1}=\sigma_{j+1}\),
lo que ocurre solo si \((i,j)=(2,n)\). En ese caso la inversión deja el
segmento \((\sigma_2,\ldots,\sigma_n)\) al revés, es decir, produce la
misma ruta recorrida en sentido contrario. Así, toda inversión
\(\rho_{ij}\) con \((i,j)\ne(2,n)\) es un movimiento 2-opt.

Sean \(x,y\in S_n\) con representaciones
\((1,\sigma_2,\ldots,\sigma_n)\) y \((1,\tau_2,\ldots,\tau_n)\). Para
\(k=2,\ldots,n-1\), si \(\sigma_k\ne\tau_k\), sea \(j>k\) la posición
actual de \(\tau_k\) y aplíquese \(\rho_{kj}\), que lleva \(\tau_k\) a
la posición \(k\) sin alterar las posiciones \(2,\ldots,k-1\). Si
\((k,j)=(2,n)\), la inversión solo cambia el sentido de recorrido y no
la ruta, así que el paso no cuenta como movimiento. Después de procesar
\(k=n-1\), la posición \(n\) queda determinada, y la sucesión obtenida
es la de \(y\). Se usaron a lo sumo \(n-2\) movimientos 2-opt, cada uno
de probabilidad positiva bajo \(R\), luego \(R^t(x,y)>0\) para algún
\(t\le n-2\).
\end{proof}

\begin{observacion}[El papel de los movimientos nulos]
\label{obs:bipartito}
La irreducibilidad no garantiza aperiodicidad. Se comprobó de manera
computacional que, para \(n=5\), el grafo de vecindad 2-opt sobre las
\(12\) rutas es bipartito; para \(n=4,6,7\) no lo es. Por lo tanto, una
propuesta uniforme sobre \(N(x)\), sin movimientos nulos, daría para
\(n=5\) una cadena de período \(2\). En la propuesta
\eqref{eq:propuesta} esto no ocurre, porque \(R(x,x)>0\).
\end{observacion}

% ------------------------------------------------
\subsection{La cadena de Metropolis a temperatura fija}
\label{subsec:metropolis}

Fijada una \emph{temperatura} \(T>0\), la regla de Metropolis
\citep{metropolis1953,hastings1970} acepta un movimiento propuesto de
\(x\) a \(y\) con probabilidad
\begin{equation}
\label{eq:aceptacion}
A_T(x,y)=\min\Bigl\{1,\;e^{-(f(y)-f(x))/T}\Bigr\}.
\end{equation}
Un movimiento que no empeora el costo se acepta siempre; uno que lo
empeora en \(\Delta>0\) se acepta con probabilidad \(e^{-\Delta/T}\).

\begin{definicion}[Cadena de Metropolis]
\label{def:metropolis}
La cadena de Metropolis a temperatura \(T\) sobre \(S_n\) tiene matriz
de transición
\begin{equation}
\label{eq:metropolis}
P_T(x,y)=
\begin{cases}
R(x,y)\,A_T(x,y), & y\ne x,\\[0.5ex]
1-\displaystyle\sum_{z\ne x}R(x,z)\,A_T(x,z), & y=x.
\end{cases}
\end{equation}
Es decir, en cada paso se propone un movimiento con la matriz \(R\); si
es nulo, la cadena permanece en \(x\), y si propone \(y\in N(x)\), la
cadena se mueve a \(y\) con probabilidad \(A_T(x,y)\) y, en caso
contrario, permanece en \(x\).
\end{definicion}

\begin{definicion}[Distribución de Boltzmann]
\label{def:boltzmann}
La \emph{distribución de Boltzmann} (o de Gibbs) a temperatura \(T\) es
\begin{equation}
\label{eq:boltzmann}
\pi_T(x)=\frac{e^{-f(x)/T}}{Z_T},
\qquad
Z_T=\sum_{z\in S_n}e^{-f(z)/T}.
\end{equation}
\end{definicion}

\begin{proposicion}[Balance detallado]
\label{prop:balance}
Para todo \(T>0\), la cadena \(P_T\) es reversible respecto de
\(\pi_T\). En consecuencia, \(\pi_T\) es estacionaria para \(P_T\).
\end{proposicion}

\begin{proof}
Basta verificar \eqref{eq:balance} para \(x\ne y\) con
\(y\in N(x)\); si \(y\notin N(x)\), ambos lados valen cero por la
proposición~\ref{prop:vecindad}(b). Supongamos, sin pérdida de
generalidad, \(f(y)\ge f(x)\). Entonces \(A_T(x,y)=e^{-(f(y)-f(x))/T}\)
y \(A_T(y,x)=1\), y usando \(R(x,y)=R(y,x)\),
\[
\pi_T(x)P_T(x,y)
=\frac{e^{-f(x)/T}}{Z_T}\,R(x,y)\,e^{-(f(y)-f(x))/T}
=\frac{e^{-f(y)/T}}{Z_T}\,R(y,x)
=\pi_T(y)P_T(y,x).
\]
La estacionariedad se sigue de la observación posterior a la
definición~\ref{def:propiedades}.
\end{proof}

Notemos que \(\pi_T\) se conoce salvo la constante \(Z_T\), cuyo
cálculo exige recorrer \(S_n\). La cadena de Metropolis solo usa
cocientes \(\pi_T(y)/\pi_T(x)=e^{-(f(y)-f(x))/T}\), en los que \(Z_T\)
se cancela; esta es la razón por la cual el método es aplicable cuando
\(|S_n|\) es astronómico.

\begin{proposicion}[Irreducibilidad y aperiodicidad]
\label{prop:aperiodica}
Para todo \(T>0\), la cadena \(P_T\) es irreducible y aperiódica.
Además, \(P_T(x,x)\ge \delta\) para todo \(x\), con
\(\delta=2/(n-1)\).
\end{proposicion}

\begin{proof}
Para \(y\in N(x)\) se tiene \(P_T(x,y)=R(x,y)A_T(x,y)>0\), pues
\(A_T>0\). Así, todo camino de movimientos 2-opt tiene probabilidad
positiva bajo \(P_T\), y la irreducibilidad se sigue de la
proposición~\ref{prop:irreducible}. Para la aperiodicidad, por
\eqref{eq:metropolis} se tiene \(P_T(x,x)\ge R(x,x)=2/(n-1)>0\), y un
estado con lazo basta para una cadena irreducible.
\end{proof}

\begin{observacion}
\label{obs:rechazos}
Aun sin movimientos nulos, la cadena de Metropolis sería aperiódica
siempre que \(f\) no sea constante: como el grafo de vecindad es conexo,
existe una arista \(\{x,y\}\) con \(f(y)>f(x)\), y entonces
\(P_T(x,x)\ge R(x,y)\bigl(1-e^{-(f(y)-f(x))/T}\bigr)>0\). Este es el
mecanismo que opera en la variante QUBO de la
sección~\ref{subsec:qubo}.
\end{observacion}

\begin{corolario}
\label{cor:metropolis}
Para todo \(T>0\), la cadena \(P_T\) tiene a
\(\pi_T\) como única distribución estacionaria,
\(P_T^t(x,\cdot)\to\pi_T\) para toda \(x\), y las frecuencias empíricas
de visita convergen casi seguramente a \(\pi_T\).
\end{corolario}

\begin{proof}
Es consecuencia directa de las proposiciones~\ref{prop:balance}
y~\ref{prop:aperiodica} y de los
teoremas~\ref{teo:unicidad}--\ref{teo:ergodico}.
\end{proof}

% ------------------------------------------------
\subsection{El papel de la temperatura}
\label{subsec:temperatura}

La temperatura interpola entre una exploración uniforme y una
concentración en los óptimos.

\begin{proposicion}
\label{prop:temperatura}
Sea \(U\) la distribución uniforme sobre \(S_n\) y \(U^\ast\) la
uniforme sobre \(\Sopt\). Entonces:
\begin{enumerate}[label=(\alph*)]
    \item \(\pi_T\to U\) cuando \(T\to\infty\);
    \item \(\pi_T\to U^\ast\) cuando \(T\to 0^+\);
    \item el costo medio \(\E_{\pi_T}[f]=\sum_x f(x)\pi_T(x)\) es no
    decreciente en \(T\), pues
    \begin{equation}
    \label{eq:derivada}
    \frac{d}{dT}\,\E_{\pi_T}[f]=\frac{\Var_{\pi_T}(f)}{T^2}\ge 0.
    \end{equation}
\end{enumerate}
\end{proposicion}

\begin{proof}
(a) Cada término \(e^{-f(x)/T}\to 1\), luego \(\pi_T(x)\to1/|S_n|\).

(b) Multiplicando numerador y denominador por \(e^{f^\ast/T}\),
\[
\pi_T(x)=\frac{e^{-(f(x)-f^\ast)/T}}{|\Sopt|+\sum_{z\notin\Sopt}e^{-(f(z)-f^\ast)/T}}.
\]
Si \(z\notin\Sopt\), entonces \(f(z)-f^\ast>0\) y
\(e^{-(f(z)-f^\ast)/T}\to 0\). Por lo tanto \(\pi_T(x)\to1/|\Sopt|\) si
\(x\in\Sopt\) y \(\pi_T(x)\to 0\) en otro caso.

(c) Sea \(\beta=1/T\) y
\(E(\beta)=Z^{-1}\sum_x f(x)e^{-\beta f(x)}\). Derivando,
\[
E'(\beta)=-\frac{\sum_x f(x)^2e^{-\beta f(x)}}{Z}
+\Bigl(\frac{\sum_x f(x)e^{-\beta f(x)}}{Z}\Bigr)^2
=-\Var_{\pi_T}(f).
\]
Como \(d\beta/dT=-1/T^2\), la regla de la cadena da
\eqref{eq:derivada}.
\end{proof}

La proposición~\ref{prop:temperatura} motiva el recocido simulado: a
temperatura baja, una muestra de \(\pi_T\) es casi siempre una ruta
óptima. La dificultad es que, como se ve a continuación, a temperatura
baja la cadena tarda exponencialmente en alcanzar el equilibrio.

% ------------------------------------------------
\subsection{Velocidad de convergencia y brecha espectral}
\label{subsec:brecha}

Como \(P_T\) es reversible respecto de \(\pi_T\), es autoadjunta en el
espacio \(\ell^2(\pi_T)\) con producto interno
\(\langle g,h\rangle_{\pi_T}=\sum_x g(x)h(x)\pi_T(x)\). Por lo tanto sus
valores propios son reales y, si \(N=|S_n|\), pueden ordenarse como
\[
1=\lambda_1>\lambda_2\ge\cdots\ge\lambda_N>-1,
\]
donde la desigualdad estricta \(\lambda_1>\lambda_2\) se debe a la
irreducibilidad y \(\lambda_N>-1\) a la aperiodicidad
\citep[cap.~12]{levin2017}.

\begin{definicion}[Brecha espectral]
\label{def:brecha}
La \emph{brecha espectral} de \(P_T\) es
\(\gamma(T)=1-\lambda_2\), y la \emph{brecha absoluta} es
\(\gamma_\ast(T)=1-\lambda_\ast\), con
\(\lambda_\ast=\max\{|\lambda_2|,|\lambda_N|\}\). El \emph{tiempo de
relajación} es \(t_{\mathrm{rel}}=1/\gamma_\ast\).
\end{definicion}

\begin{proposicion}[Espectro no muy negativo]
\label{prop:espectro_lazy}
Todo valor propio de \(P_T\) cumple \(\lambda\ge 2\delta-1\), con
\(\delta=2/(n-1)\). En particular, para \(n=5\) todos los valores propios
son no negativos y \(\gamma_\ast(T)=\gamma(T)\).
\end{proposicion}

\begin{proof}
Por la proposición~\ref{prop:aperiodica}, la matriz
\(P'=(P_T-\delta I)/(1-\delta)\) tiene entradas no negativas y filas que
suman \(1\), luego es estocástica y sus valores propios tienen módulo
a lo sumo \(1\). Como \(P_T=\delta I+(1-\delta)P'\), cada valor propio de
\(P_T\) es de la forma \(\delta+(1-\delta)\mu\) con \(\mu\ge-1\), y por
lo tanto es al menos \(2\delta-1\). Para \(n=5\), \(\delta=1/2\).
\end{proof}

\begin{proposicion}[Cota espectral]
\label{prop:cota_espectral}
Sea \(\pi_{\min}=\min_x\pi_T(x)\). Para todo \(t\ge 0\),
\[
\max_{x\in S_n}\bigl\|P_T^t(x,\cdot)-\pi_T\bigr\|_{\TV}
\le\frac{\lambda_\ast^{\,t}}{2\,\pi_{\min}}
\le\frac{e^{-\gamma_\ast t}}{2\,\pi_{\min}}.
\]
\end{proposicion}

\begin{proof}
Sea \(\{\varphi_j\}_{j=1}^N\) una base ortonormal de \(\ell^2(\pi_T)\)
formada por vectores propios de \(P_T\), con \(\varphi_1\equiv1\). La
descomposición espectral da
\[
\frac{P_T^t(x,y)}{\pi_T(y)}=1+\sum_{j=2}^N\varphi_j(x)\varphi_j(y)\lambda_j^t.
\]
Aplicando la ortonormalidad a la función \(\mathbf 1_{\{x\}}\) se
obtiene \(\sum_{j}\varphi_j(x)^2=1/\pi_T(x)\). Por la desigualdad de
Cauchy--Schwarz,
\[
\Bigl|\frac{P_T^t(x,y)}{\pi_T(y)}-1\Bigr|
\le\lambda_\ast^{\,t}\,\frac{1}{\sqrt{\pi_T(x)\pi_T(y)}}
\le\frac{\lambda_\ast^{\,t}}{\pi_{\min}}.
\]
Multiplicando por \(\pi_T(y)/2\) y sumando sobre \(y\) se obtiene la
primera desigualdad; la segunda se sigue de
\(1-\gamma_\ast\le e^{-\gamma_\ast}\). Véase también
\citet[cap.~12]{levin2017}.
\end{proof}

Así, el número de pasos necesarios para que la cadena olvide su estado
inicial es del orden de \(t_{\mathrm{rel}}\), salvo un factor
logarítmico en \(1/\pi_{\min}\). A temperatura baja, \citet{ingrassia1994}
muestra que para la cadena de Metropolis la brecha que domina es
\(\gamma=1-\lambda_2\). Su comportamiento depende de la geometría del
paisaje de costos, que se describe con las siguientes nociones
\citep{hajek1988}.

\begin{definicion}[Altura de comunicación y profundidad]
\label{def:profundidad}
Para un camino \(x=x_0,x_1,\ldots,x_k=y\) con \(x_{i+1}\in N(x_i)\), su
\emph{altura} es \(\max_i f(x_i)\). La \emph{altura de comunicación}
\(H(x,y)\) es la mínima altura entre todos los caminos de \(x\) a
\(y\). Un estado \(x\) es un \emph{mínimo local} si \(f(x)\le f(y)\)
para todo \(y\in N(x)\). Para \(x\notin\Sopt\), su \emph{profundidad} es
\[
h(x)=\min_{y:\,f(y)<f(x)}H(x,y)-f(x),
\]
es decir, la mínima subida necesaria para escapar de \(x\) hacia un
estado estrictamente mejor. La \emph{profundidad crítica} es
\[
h^\ast=\max_{x\notin\Sopt}h(x),
\]
la profundidad del pozo no óptimo más hondo.
\citet{hajek1988} denota esta constante por \(d^\ast\).
\end{definicion}

Si \(x\) no es mínimo local, existe un vecino mejor y \(h(x)=0\). Por
lo tanto, \(h^\ast>0\) si y solo si existe un mínimo local que no es
global. La profundidad crítica gobierna la brecha espectral a
temperatura baja.

\begin{teorema}[Brecha a temperatura baja]
\label{teo:brecha_baja}
Supongamos que \(f\) tiene un único mínimo global. Entonces
\[
\lim_{T\to 0^+}T\log\gamma(T)=-h^\ast,
\]
es decir, \(\gamma(T)=e^{-h^\ast/T+o(1/T)}\) y
\(t_{\mathrm{rel}}\approx e^{h^\ast/T}\).
\end{teorema}

Este resultado se debe a \citet{holley1988} y \citet{holley1989}; la
formulación y la identificación de la constante con la profundidad
crítica de Hajek siguen a \citet{bertsimas1993}. El
teorema~\ref{teo:brecha_baja} da una predicción cuantitativa que se
verifica en la sección~\ref{sec:resultados}: en una gráfica de
\(-\log\gamma(T)\) frente a \(1/T\), la pendiente a temperatura baja
debe aproximarse a \(h^\ast\), que para \(n=5\) puede calcularse de
forma exacta recorriendo el grafo de vecindad.

% ------------------------------------------------
\subsection{El recocido simulado como cadena no homogénea}
\label{subsec:recocido}

El teorema~\ref{teo:brecha_baja} expone el dilema: una temperatura baja
concentra \(\pi_T\) en \(\Sopt\), pero hace que la cadena tarde del
orden de \(e^{h^\ast/T}\) pasos en equilibrarse. El \emph{recocido
simulado} \citep{kirkpatrick1983,cerny1985} resuelve el dilema
disminuyendo la temperatura de forma gradual.

\begin{definicion}[Recocido simulado]
\label{def:recocido}
Sea \((T_k)_{k\ge0}\) una sucesión de temperaturas positivas, no
creciente y con \(T_k\to0\), llamada \emph{esquema de enfriamiento}. El
\emph{recocido simulado} es la cadena de Markov \emph{no homogénea}
\((X_k)_{k\ge 0}\) con
\[
\Prob(X_{k+1}=y\mid X_k=x)=P_{T_k}(x,y),
\]
donde \(P_{T_k}\) es la matriz de Metropolis \eqref{eq:metropolis} a
temperatura \(T_k\).
\end{definicion}

Como las matrices de transición cambian con \(k\), los
teoremas~\ref{teo:unicidad}--\ref{teo:ergodico} ya no se aplican. La
pregunta natural es si \(\Prob(X_k\in\Sopt)\to1\). La respuesta
completa la da el teorema de Hajek; sus hipótesis de irreducibilidad y
de reversibilidad débil se cumplen aquí por las
proposiciones~\ref{prop:vecindad}(b) y~\ref{prop:irreducible}.

\begin{teorema}[\citealp{hajek1988}]
\label{teo:hajek}
Sea \((T_k)\) no creciente con \(T_k\to0\). Entonces
\[
\lim_{k\to\infty}\Prob(X_k\in\Sopt)=1
\quad\Longleftrightarrow\quad
\sum_{k=0}^{\infty}e^{-h^\ast/T_k}=\infty.
\]
\end{teorema}

\begin{corolario}[Enfriamiento logarítmico frente a geométrico]
\label{cor:enfriamiento}
Supongamos \(h^\ast>0\).
\begin{enumerate}[label=(\alph*)]
    \item Con el esquema logarítmico \(T_k=c/\log(k+2)\), el recocido
    converge a \(\Sopt\) si y solo si \(c\ge h^\ast\).
    \item Con el esquema geométrico \(T_k=T_0\alpha^k\), con
    \(0<\alpha<1\), no se cumple \(\Prob(X_k\in\Sopt)\to1\).
\end{enumerate}
\end{corolario}

\begin{proof}
(a) Se tiene \(e^{-h^\ast/T_k}=(k+2)^{-h^\ast/c}\), y la serie
\(\sum_k (k+2)^{-h^\ast/c}\) diverge si y solo si \(h^\ast/c\le1\).

(b) Se tiene \(e^{-h^\ast/T_k}=\exp(-h^\ast\alpha^{-k}/T_0)\). Como
\(\alpha^{-k}\ge k\) para \(k\) grande, el término general está acotado
por \(e^{-(h^\ast/T_0)k}\), cuya serie converge. Por el
teorema~\ref{teo:hajek}, no hay convergencia a \(\Sopt\) con
probabilidad tendiente a uno.
\end{proof}

El corolario~\ref{cor:enfriamiento} muestra una tensión entre teoría y
práctica. El esquema logarítmico garantiza convergencia pero es
extremadamente lento: para bajar de \(T_0\) a \(T_0/2\) hay que elevar
al cuadrado el número de pasos. Por eso, en la práctica se usa el
esquema geométrico, que carece de garantía asintótica pero, con un
presupuesto finito de pasos, suele dar soluciones de muy buena calidad
\citep{delahaye2019}. Una forma habitual de implementarlo es mantener
cada temperatura durante \(L\) pasos antes de multiplicarla por
\(\alpha\), para que la cadena se aproxime a \(\pi_{T}\) en cada
nivel.

% ------------------------------------------------
\subsection{Variante sobre el modelo QUBO}
\label{subsec:qubo}

El recocido también puede aplicarse a una formulación del TSP como
problema de optimización cuadrática binaria sin restricciones (QUBO),
que es la forma de entrada de los recocedores cuánticos
\citep{lucas2014,glover2022,bochkarev2026}. Esta variante cambia el
espacio de estados, y con él la geometría del paisaje.

\paragraph{Formulación.} Se fija la ciudad \(1\) en la posición \(1\) y
se usan variables binarias \(y_{ik}\in\{0,1\}\), para
\(i,k\in\{2,\ldots,n\}\), con \(y_{ik}=1\) si la ciudad \(i\) ocupa la
posición \(k\). El vector \(y\) tiene \(m=(n-1)^2\) componentes. El
costo de recorrido es
\begin{equation}
\label{eq:costo_qap}
f_{\mathrm Q}(y)=\sum_{i=2}^n c_{1i}\,y_{i2}
+\sum_{i=2}^n c_{i1}\,y_{in}
+\sum_{k=2}^{n-1}\sum_{i\ne j}c_{ij}\,y_{ik}\,y_{j,k+1},
\end{equation}
y las restricciones de asignación se incorporan como penalización:
\begin{equation}
\label{eq:penalizacion}
\Phi(y)=\sum_{i=2}^n\Bigl(\sum_{k=2}^n y_{ik}-1\Bigr)^2
+\sum_{k=2}^n\Bigl(\sum_{i=2}^n y_{ik}-1\Bigr)^2,
\qquad
E_M(y)=f_{\mathrm Q}(y)+M\,\Phi(y),
\end{equation}
con \(M>0\). Usando \(y_{ik}^2=y_{ik}\), la energía se escribe como
\(E_M(y)=y^{\top}Qy+\text{const}\) con \(Q\) simétrica, que es la forma
QUBO. Sea \(F\subset\{0,1\}^m\) el conjunto de vectores
\emph{factibles}, aquellos con \(\Phi(y)=0\), que son las matrices de
permutación. Para \(y\in F\), \(f_{\mathrm Q}(y)\) es el costo de la
ruta correspondiente; cada ruta corresponde a dos elementos de \(F\),
uno por cada sentido, así que \(|F|=(n-1)!\).

\begin{proposicion}[Equivalencia para \(M\) grande]
\label{prop:equivalencia}
Si \(M>f^\ast\), todo minimizador global de \(E_M\) es factible y
corresponde a una ruta óptima.
\end{proposicion}

\begin{proof}
Si \(y\notin F\), \(\Phi(y)\) es un entero positivo, luego
\(\Phi(y)\ge1\). Como \(c_{ij}\ge0\), se tiene \(f_{\mathrm Q}\ge0\) y
\(E_M(y)\ge M>f^\ast\). En cambio, el elemento de \(F\) asociado a una
ruta óptima tiene energía \(f^\ast\). Así, ningún vector infactible es
minimizador, y entre los factibles \(E_M=f_{\mathrm Q}\) se minimiza en
las rutas óptimas.
\end{proof}

La condición \(M>f^\ast\) es suficiente, no necesaria: el umbral exacto
\(M^\ast\) a partir del cual se preserva la equivalencia puede ser
menor. Como \(f^\ast\) es desconocido, en la práctica se usa una cota
superior, como el costo de cualquier ruta conocida; en este trabajo se
toma \(M=U+1\), con \(U\) el costo de la ruta del vecino más cercano. Elegir \(M\)
demasiado pequeño hace que el mínimo de \(E_M\) sea infactible; la
proposición siguiente muestra que elegirlo demasiado grande también
tiene un costo, esta vez para la cadena.

\paragraph{La cadena.} El espacio de estados es \(\{0,1\}^m\) y la
propuesta consiste en voltear un bit elegido uniformemente al azar:
\(R_{\mathrm Q}(y,y')=1/m\) si \(y\) y \(y'\) difieren en exactamente una
coordenada. El grafo de vecindad es el hipercubo, que es conexo, y
\(R_{\mathrm Q}\) es simétrica. Por lo tanto, la cadena de Metropolis
con energía \(E_M\) es reversible respecto de
\(\pi_T(y)\propto e^{-E_M(y)/T}\). Como el hipercubo es bipartito, la
aperiodicidad proviene nuevamente de los rechazos, con el mismo
argumento de la proposición~\ref{prop:aperiodica}. La diferencia con la
cadena sobre rutas está en el paisaje.

\begin{proposicion}[Barrera de penalización]
\label{prop:barrera}
Sean \(y\in F\) y \(y'\) un vector que difiere de \(y\) en un solo bit.
Entonces
\begin{enumerate}[label=(\alph*)]
    \item \(\Phi(y')=2\) y
    \(E_M(y')-E_M(y)\ge 2M-2\cmax\);
    \item dos elementos distintos de \(F\) difieren en al menos \(4\)
    coordenadas, de modo que todo camino de volteos de un bit entre
    ellos pasa por vectores infactibles;
    \item si \(M>\cmax\), cada elemento de \(F\) es un mínimo local
    estricto de \(E_M\), y la profundidad de cada uno que no es óptimo
    es al menos \(2M-2\cmax\). En particular, si existen rutas no
    óptimas, la profundidad crítica de la cadena QUBO cumple
    \[
    h^\ast_M\ge 2M-2\cmax .
    \]
\end{enumerate}
\end{proposicion}

\begin{proof}
(a) Si se apaga un bit \(y_{ik}=1\), la fila \(i\) y la columna \(k\)
quedan con suma \(0\); si se enciende un bit \(y_{ik}=0\), ambas quedan
con suma \(2\). En los dos casos, dos términos de \(\Phi\) valen \(1\) y
los demás \(0\), luego \(\Phi(y')=2\). Encender un bit solo agrega
términos no negativos a \eqref{eq:costo_qap}. Apagar \(y_{ik}\) elimina,
a lo sumo, los dos términos que unen la ciudad \(i\) con sus vecinas en
las posiciones \(k-1\) y \(k+1\), incluidos los términos lineales si
\(k=2\) o \(k=n\). Así, \(f_{\mathrm Q}(y')\ge f_{\mathrm Q}(y)-2\cmax\)
y \(E_M(y')-E_M(y)\ge 2M-2\cmax\).

(b) Dos matrices de permutación distintas difieren en al menos dos
filas, y en cada una de esas filas el \(1\) cambia de columna, lo que
modifica dos entradas por fila.

(c) Por (a), todos los vecinos de \(y\in F\) tienen energía
estrictamente mayor cuando \(M>\cmax\). Cualquier camino que salga de
\(y\) pasa primero por un vecino, así que su altura supera \(E_M(y)\) en
al menos \(2M-2\cmax\). Si \(y\) no es óptimo, la definición de
profundidad da \(h(y)\ge 2M-2\cmax\).
\end{proof}

La proposición~\ref{prop:barrera}, junto con el
teorema~\ref{teo:brecha_baja}, explica la dificultad de la variante
QUBO. Su tiempo de relajación a temperatura baja crece al menos como
\(e^{(2M-2\cmax)/T}\). Así, aumentar \(M\) para garantizar la
equivalencia (proposición~\ref{prop:equivalencia}) hace la cadena
exponencialmente más lenta, mientras que reducirlo puede volver
infactible el mínimo. Se espera entonces que el desempeño del recocido
sobre el QUBO, como función de \(M\), tenga forma de U, con el mejor
valor cerca del umbral de equivalencia. En contraste, la cadena sobre
rutas nunca sale del conjunto factible: cada movimiento 2-opt lleva de
una ruta a otra, y las únicas barreras son las del propio costo \(f\).
Esta diferencia se contrasta experimentalmente en la
sección~\ref{sec:resultados}.

% ------------------------------------------------
% 4. METODOLOGÍA E IMPLEMENTACIÓN COMPUTACIONAL
% ------------------------------------------------

\section{Metodología e implementación computacional}
\label{sec:metodologia}

TODO: pseudocódigo, diseño de los experimentos, semilla y enlace al notebook de Colab.

% ------------------------------------------------
% 5. RESULTADOS DE LAS SIMULACIONES Y ANÁLISIS
% ------------------------------------------------

\section{Resultados de las simulaciones y análisis}
\label{sec:resultados}

TODO

% ------------------------------------------------
% 6. CONCLUSIONES
% ------------------------------------------------

\section{Conclusiones}
\label{sec:conclusiones}

TODO


\bibliographystyle{plainnat-es}
\bibliography{referencias}

\end{document}
